\documentclass[11pt,a4paper,oneside,headings=big,chapterprefix=true]{scrreprt}

% ---------- Core packages ----------
\usepackage[utf8]{inputenc}
\usepackage[T1]{fontenc}
\usepackage[left=2.6cm,right=2.6cm,top=2.8cm,bottom=2.8cm,headheight=28pt]{geometry}
\usepackage{amsmath}
\usepackage{amssymb}
\usepackage{xcolor}
\usepackage[table]{xcolor}
\usepackage{graphicx}
\usepackage{booktabs}
\usepackage{colortbl}
\usepackage{array}
\usepackage{tabularx}
\usepackage{enumitem}
\usepackage{tikz}
\usetikzlibrary{positioning,calc,shapes.geometric,decorations.pathmorphing,backgrounds,patterns}
\usepackage[most]{tcolorbox}
\usepackage{fontawesome5}
\usepackage{scrlayer-scrpage}
\usepackage{hyperref}

% ---------- Palette ----------
\definecolor{steel}{RGB}{33,69,94}
\definecolor{amberaccent}{RGB}{198,138,42}
\definecolor{fogbg}{RGB}{238,241,243}
\definecolor{inkgray}{RGB}{45,49,54}

% ---------- Typography ----------
\addtokomafont{disposition}{\color{steel}}
\addtokomafont{chapter}{\color{steel}\bfseries}
\addtokomafont{chapterprefix}{\normalfont\normalsize\color{amberaccent}\bfseries}
\setkomafont{chapterentry}{\color{steel}}
\setkomafont{caption}{\small\color{inkgray}}
\setkomafont{captionlabel}{\small\bfseries\color{steel}}
\renewcommand*{\chapterformat}{\mbox{\chapappifchapterprefix{\ }\thechapter\autodot\enskip}}

\clearpairofpagestyles
\ihead{\small\color{steel}\slshape\leftmark}
\ohead{\small\color{steel}\thepage}
\ofoot{\small\color{inkgray}\textsc{Nexa Systems} \textbar\ Atlas-7 Field Guide}
\ifoot{\small\color{inkgray}Rev.\ 3.2}
\setkomafont{pageheadfoot}{\normalfont}
\automark[section]{chapter}
\pagestyle{scrheadings}

\hypersetup{
  colorlinks=true,
  linkcolor=steel,
  urlcolor=amberaccent,
  citecolor=steel,
  pdftitle={Atlas-7 Field Telemetry Unit --- Technical Manual and Field Guide},
  pdfauthor={Nexa Systems Technical Publications}
}

% ---------- Alert boxes ----------
\newtcolorbox{dangerbox}[1][]{%
  colback=amberaccent!10,colframe=amberaccent!85!black,coltitle=white,
  fonttitle=\bfseries,boxrule=0.9pt,arc=1.5pt,left=6pt,right=6pt,top=4pt,bottom=4pt,
  title=\faExclamationTriangle\ \ DANGER,#1
}
\newtcolorbox{warnbox}[1][]{%
  colback=amberaccent!8,colframe=amberaccent,coltitle=inkgray,
  fonttitle=\bfseries,boxrule=0.8pt,arc=1.5pt,left=6pt,right=6pt,top=4pt,bottom=4pt,
  title=\faExclamationTriangle\ \ WARNING,#1
}
\newtcolorbox{notebox}[1][]{%
  colback=steel!6,colframe=steel,coltitle=white,
  fonttitle=\bfseries,boxrule=0.8pt,arc=1.5pt,left=6pt,right=6pt,top=4pt,bottom=4pt,
  title=\faInfoCircle\ \ NOTE,#1
}
\newtcolorbox{tipbox}[1][]{%
  colback=fogbg,colframe=steel!60,coltitle=steel,
  fonttitle=\bfseries,boxrule=0.7pt,arc=1.5pt,left=6pt,right=6pt,top=4pt,bottom=4pt,
  title=\faWrench\ \ FIELD TIP,#1
}

\setlist[itemize,1]{label=\textcolor{steel}{\small\faCheckCircle},leftmargin=1.6em}
\setlist[itemize,2]{label=\textcolor{amberaccent}{\textbullet},leftmargin=1.4em}
\setlist[enumerate,1]{label=\textcolor{steel}{\bfseries\arabic*.},leftmargin=1.8em}

\newcommand{\revtable}[4]{#1 & #2 & #3 & #4 \\}

\begin{document}

% ============================================================
% TITLE PAGE
% ============================================================
\begin{titlepage}
\begin{tikzpicture}[remember picture,overlay]
  \fill[steel] (current page.north west) rectangle ([yshift=-4.2cm]current page.north east);
  \fill[amberaccent] ([yshift=-4.2cm]current page.north west) rectangle ([yshift=-4.32cm]current page.north east);
  \node[anchor=north west,text=white,font=\bfseries\Large] at ([xshift=2.4cm,yshift=-1.5cm]current page.north west) {NEXA SYSTEMS};
  \node[anchor=north west,text=white!85!steel,font=\small] at ([xshift=2.4cm,yshift=-2.5cm]current page.north west) {Technical Publications Division};
\end{tikzpicture}

\vspace*{5.6cm}

\begin{center}
{\color{steel}\rule{0.55\textwidth}{1.4pt}}\\[0.6cm]
{\fontsize{30}{34}\selectfont\bfseries\color{steel} Atlas-7}\\[0.15cm]
{\Large\color{inkgray} Field Telemetry Unit}\\[0.5cm]
{\large\color{amberaccent}\bfseries TECHNICAL MANUAL \& FIELD GUIDE}\\[0.6cm]
{\color{steel}\rule{0.55\textwidth}{1.4pt}}
\end{center}

\vspace{1.6cm}
\begin{center}
\begin{tikzpicture}
  \node[draw=steel,rounded corners=2pt,inner sep=10pt,text width=9.4cm,align=left,fill=fogbg] {%
    \small
    \begin{tabular}{@{}ll@{}}
      \textbf{\color{steel}Document Number:} & NX-ATL7-TM-032 \\
      \textbf{\color{steel}Revision:} & 3.2 \\
      \textbf{\color{steel}Issue Date:} & 14 March 2026 \\
      \textbf{\color{steel}Applicable Hardware:} & Atlas-7 Rev.\ C and later \\
      \textbf{\color{steel}Classification:} & Customer Field Documentation \\
    \end{tabular}
  };
\end{tikzpicture}
\end{center}

\vfill
\begin{center}
{\small\color{inkgray} Prepared by the Nexa Systems Technical Publications Team\\
Reviewed by Elena Vasquez, Lead Field Engineer, Meridian Deployment Group\\[0.3cm]
\copyright\ Nexa Systems. All rights reserved. Unauthorized reproduction prohibited.}
\end{center}
\vspace{1cm}
\end{titlepage}

% ============================================================
% REVISION HISTORY
% ============================================================
\chapter*{Revision History}
\addcontentsline{toc}{chapter}{Revision History}
\begin{center}
\renewcommand{\arraystretch}{1.3}
\begin{tabularx}{\textwidth}{@{}l l l X@{}}
\toprule
\rowcolor{steel}
\textcolor{white}{\bfseries Rev.} & \textcolor{white}{\bfseries Date} & \textcolor{white}{\bfseries Author} & \textcolor{white}{\bfseries Summary of Changes} \\
\midrule
\revtable{1.0}{02 Jun 2025}{J.\ Okafor}{Initial release for Atlas-7 Rev.\ A hardware.}
\revtable{2.0}{19 Nov 2025}{J.\ Okafor}{Added LoRaWAN provisioning chapter; revised specifications.}
\revtable{3.0}{08 Feb 2026}{E.\ Vasquez}{New troubleshooting matrix; updated safety notices.}
\rowcolor{fogbg}
\revtable{3.2}{14 Mar 2026}{E.\ Vasquez}{Corrected mounting torque values; added quick reference card.}
\bottomrule
\end{tabularx}
\end{center}

% ============================================================
\tableofcontents

% ============================================================
\chapter{About This Guide}
\label{ch:about}

This manual documents the installation, operation, and maintenance of the \textbf{Nexa Atlas-7 Field Telemetry Unit}, a ruggedized environmental and asset-monitoring device intended for unattended outdoor deployment. It is written for field technicians, systems integrators, and site operators who install or service Atlas-7 units in the course of their work.

\section{Who Should Use This Manual}

\begin{itemize}
  \item \textbf{Field technicians} performing initial installation, commissioning, or scheduled maintenance.
  \item \textbf{Systems integrators} incorporating the Atlas-7 into a larger sensor network or SCADA environment.
  \item \textbf{Site operators} responsible for day-to-day monitoring and first-line troubleshooting.
\end{itemize}

No specialized electronics training is required to complete standard installation procedures; however, Chapter~\ref{ch:maintenance} assumes familiarity with basic hand tools and low-voltage DC wiring practice.

\section{How This Manual Is Organized}

\begin{center}
\renewcommand{\arraystretch}{1.25}
\begin{tabularx}{\textwidth}{@{}l X@{}}
\toprule
\rowcolor{steel}
\textcolor{white}{\bfseries Chapter} & \textcolor{white}{\bfseries Contents} \\
\midrule
2 & System overview, key features, and hardware architecture \\
\rowcolor{fogbg}
3 & Pre-installation planning, mounting, and network setup \\
4 & Day-to-day operation, control indicators, and operating modes \\
\rowcolor{fogbg}
5 & Maintenance schedule and troubleshooting reference \\
App.~A & Full technical specifications \\
\rowcolor{fogbg}
App.~B & Quick reference card (tear-out summary) \\
\bottomrule
\end{tabularx}
\end{center}

\section{Notice Conventions}

Throughout this guide, callouts flag information that affects safety, data integrity, or successful commissioning. Each type is used consistently:

\begin{dangerbox}
Indicates a hazard that \textbf{will} result in serious injury or permanent equipment damage if ignored. Danger notices always precede the step they apply to, never follow it.
\end{dangerbox}

\begin{warnbox}
Indicates a hazard or action that \textbf{may} result in injury, data loss, or equipment damage if precautions are not observed.
\end{warnbox}

\begin{notebox}
Supplementary information that clarifies a procedure but carries no safety implication.
\end{notebox}

\begin{tipbox}
Practical, field-tested advice that speeds up a task or avoids a common mistake.
\end{tipbox}

% ============================================================
\chapter{System Overview}
\label{ch:overview}

\section{Product Description}

The Atlas-7 is a self-contained, solar-assisted telemetry node designed for multi-year unattended operation in remote or harsh environments. It combines a modular sensor bay, an industrial communications stack, and a weatherproof enclosure rated to IP66. A single Atlas-7 unit can host up to six interchangeable sensor modules and report readings over cellular LTE-M, LoRaWAN, or a wired RS-485 backhaul, depending on the communications card fitted.

\subsection*{Key Features}
\begin{itemize}
  \item Modular sensor bay accepting up to six field-swappable probes (temperature, humidity, soil moisture, vibration, water level, or custom analog input).
  \item Dual-radio communications card: LTE-M/NB-IoT primary link with LoRaWAN failover.
  \item Solar trickle-charge system rated for 90 days of full function with zero sunlight.
  \item Tamper-evident enclosure latch with optional tilt and open-door alarms.
  \item Over-the-air firmware updates via the Nexa Cloud fleet console.
\end{itemize}

\section{Hardware Architecture}

Figure~\ref{fig:architecture} shows the functional blocks of a fully provisioned Atlas-7 unit and how it exchanges data with Nexa Cloud.

\begin{figure}[h]
\centering
\begin{tikzpicture}[
  node distance=10mm and 14mm,
  block/.style={draw=steel,rounded corners=3pt,fill=fogbg,minimum height=1cm,minimum width=3.1cm,align=center,font=\small,thick},
  core/.style={draw=steel,rounded corners=3pt,fill=steel,text=white,minimum height=1.2cm,minimum width=3.3cm,align=center,font=\small\bfseries,thick},
  cloud/.style={draw=amberaccent,rounded corners=3pt,fill=amberaccent!12,minimum height=1.1cm,minimum width=3.3cm,align=center,font=\small,thick}
]
  \node[block] (sensors) {Sensor Bay\\(up to 6 probes)};
  \node[core, right=of sensors] (core) {Atlas-7\\Core Unit};
  \node[block, above right=of core] (comms) {Comms Card\\LTE-M / LoRaWAN};
  \node[block, below right=of core] (power) {Power \& Solar\\Management};
  \node[cloud, right=24mm of comms] (cloud) {Nexa Cloud\\Fleet Console};

  \draw[->,thick,steel] (sensors) -- (core);
  \draw[->,thick,steel] (core) -- (comms);
  \draw[->,thick,steel] (core) -- (power);
  \draw[->,thick,amberaccent] (comms) -- (cloud);
\end{tikzpicture}
\caption{Atlas-7 functional block diagram.}
\label{fig:architecture}
\end{figure}

\section{Specifications Summary}

\begin{center}
\renewcommand{\arraystretch}{1.25}
\begin{tabularx}{\textwidth}{@{}l X@{}}
\toprule
\rowcolor{steel}
\textcolor{white}{\bfseries Parameter} & \textcolor{white}{\bfseries Value} \\
\midrule
Enclosure rating & IP66, UV-stabilized polycarbonate \\
\rowcolor{fogbg}
Operating temperature & $-30\,^{\circ}\mathrm{C}$ to $+60\,^{\circ}\mathrm{C}$ \\
Power input & 12\,V DC nominal, integrated 8\,W solar panel \\
\rowcolor{fogbg}
Battery & 18\,Ah LiFePO4, field-replaceable \\
Sensor bay & 6 ports, screw-terminal, hot-swappable \\
\rowcolor{fogbg}
Primary radio & LTE-M / NB-IoT, 3GPP Release 13 \\
Failover radio & LoRaWAN 1.0.3, EU868 / US915 selectable \\
\rowcolor{fogbg}
Dimensions & 340\,mm $\times$ 220\,mm $\times$ 118\,mm \\
Weight (unpopulated) & 2.4\,kg \\
\bottomrule
\end{tabularx}
\end{center}

A complete specification listing, including tolerances and certification marks, appears in Appendix~\ref{app:specs}.

% ============================================================
\chapter{Installation and Setup}
\label{ch:install}

\section{Pre-Installation Checklist}

Before travelling to a deployment site, confirm the following:

\begin{itemize}
  \item Mounting location provides a clear southern sky view for the solar panel (northern sky view in the southern hemisphere).
  \item Cellular coverage or a LoRaWAN gateway is confirmed within range using the Nexa Cloud coverage map.
  \item The unit has been pre-registered in Nexa Cloud with its unique device ID (printed on the enclosure label).
  \item A torque driver (2--6\,N\textperiodcentered m) and M8 stainless mounting hardware are on hand.
\end{itemize}

\begin{warnbox}
Do not power on the Atlas-7 indoors under fluorescent lighting for extended periods. The ambient-light sensor may falsely trigger daylight calibration, producing a skewed offset for the first 24 hours of field operation.
\end{warnbox}

\section{Mounting Procedure}

\begin{enumerate}
  \item Select a pole or wall location between 1.5\,m and 3\,m above grade, clear of overhanging vegetation.
  \item Attach the mounting bracket using the supplied M8 bolts, torqued to 4.5\,N\textperiodcentered m.
  \item Orient the enclosure so the solar panel faces true south ($\pm 15^{\circ}$) and is tilted to the site latitude.
  \item Insert sensor probes into the desired bay ports and hand-tighten the compression glands.
  \item Close and latch the enclosure door; confirm the tamper switch reads \emph{closed} on the status LED.
\end{enumerate}

\begin{dangerbox}
Never mount the Atlas-7 within 3\,m of overhead power lines. The mounting bracket is conductive stainless steel and provides no electrical isolation.
\end{dangerbox}

\section{Network Provisioning}

\begin{center}
\renewcommand{\arraystretch}{1.25}
\begin{tabularx}{\textwidth}{@{}l X@{}}
\toprule
\rowcolor{steel}
\textcolor{white}{\bfseries Step} & \textcolor{white}{\bfseries Action} \\
\midrule
1 & Power on the unit; hold the pairing button for 3 seconds until the status LED pulses blue. \\
\rowcolor{fogbg}
2 & In Nexa Cloud, open \textbf{Fleet~\textgreater~Add Device} and scan the QR code on the enclosure label. \\
3 & Select the communications profile: \texttt{LTE-PRIMARY}, \texttt{LORA-ONLY}, or \texttt{DUAL-FAILOVER}. \\
\rowcolor{fogbg}
4 & Confirm first telemetry packet arrives in the console (typically within 90 seconds). \\
5 & Lock the configuration to prevent accidental profile changes from the local button. \\
\bottomrule
\end{tabularx}
\end{center}

\begin{tipbox}
If first telemetry does not arrive within five minutes, check signal strength on the status LED (see Table~\ref{tab:leds} in Chapter~\ref{ch:operation}) before assuming a network fault --- most delayed first-packets are solar-panel shading issues, not radio problems.
\end{tipbox}

% ============================================================
\chapter{Operation}
\label{ch:operation}

\section{Startup Sequence}

On power-up, the Atlas-7 performs a self-test lasting approximately 45 seconds: sensor bay enumeration, battery health check, and a radio handshake attempt. The status LED reports progress as summarized in Table~\ref{tab:leds}.

\begin{table}[h]
\centering
\renewcommand{\arraystretch}{1.3}
\begin{tabularx}{\textwidth}{@{}l l X@{}}
\toprule
\rowcolor{steel}
\textcolor{white}{\bfseries LED State} & \textcolor{white}{\bfseries Color} & \textcolor{white}{\bfseries Meaning} \\
\midrule
Solid & \textcolor{amberaccent!80!black}{\bfseries Amber} & Self-test in progress \\
\rowcolor{fogbg}
Slow pulse & \textcolor{steel}{\bfseries Blue} & Pairing mode, awaiting Nexa Cloud handshake \\
Solid & \textcolor{steel}{\bfseries Green} & Normal operation, strong signal \\
\rowcolor{fogbg}
Slow pulse & \textcolor{steel}{\bfseries Green} & Normal operation, marginal signal \\
Fast pulse & \textcolor{amberaccent!80!black}{\bfseries Amber} & Sensor fault detected, telemetry continues \\
\rowcolor{fogbg}
Solid & Red & Battery critical or radio failure \\
\bottomrule
\end{tabularx}
\caption{Status LED reference.}
\label{tab:leds}
\end{table}

\section{Operating Modes}

\begin{itemize}
  \item \textbf{Continuous mode} --- reports all sensor channels every 5 minutes. Highest power draw; recommended only where solar exposure exceeds 4 peak-sun-hours daily.
  \item \textbf{Scheduled mode} --- reports at configurable intervals (default 30 minutes), with event-triggered reports for threshold breaches.
  \item \textbf{Low-power mode} --- reports every 4 hours and disables the failover radio. Automatically engaged when battery state of charge falls below 25\%.
\end{itemize}

\begin{notebox}
Mode transitions are logged to Nexa Cloud with a timestamp and triggering condition, allowing post-deployment review of power events without a site visit.
\end{notebox}

% ============================================================
\chapter{Maintenance and Troubleshooting}
\label{ch:maintenance}

\section{Scheduled Maintenance}

\begin{center}
\renewcommand{\arraystretch}{1.3}
\begin{tabularx}{\textwidth}{@{}l l X@{}}
\toprule
\rowcolor{steel}
\textcolor{white}{\bfseries Interval} & \textcolor{white}{\bfseries Task} & \textcolor{white}{\bfseries Notes} \\
\midrule
Monthly & Visual inspection & Check for debris on solar panel, gland integrity \\
\rowcolor{fogbg}
Quarterly & Sensor probe cleaning & Use soft brush; never use solvents on probe housings \\
Annually & Battery capacity test & Via Nexa Cloud \textbf{Diagnostics} tab; replace below 80\% rated capacity \\
\rowcolor{fogbg}
Annually & Enclosure gasket check & Replace if compression set exceeds 30\% \\
\bottomrule
\end{tabularx}
\end{center}

\section{Troubleshooting Reference}

\begin{center}
\renewcommand{\arraystretch}{1.35}
\begin{tabularx}{\textwidth}{@{}X X X@{}}
\toprule
\rowcolor{steel}
\textcolor{white}{\bfseries Symptom} & \textcolor{white}{\bfseries Probable Cause} & \textcolor{white}{\bfseries Recommended Action} \\
\midrule
No telemetry, LED off & Battery fully discharged & Verify solar panel is unshaded; check panel output with a multimeter ($>$14\,V DC in daylight expected) \\
\rowcolor{fogbg}
LED red, radio failure & SIM or LoRaWAN key mismatch & Re-run provisioning in Section~\ref{ch:install}, Network Provisioning \\
Sensor reads \texttt{NULL} & Probe not seated or compression gland loose & Reseat probe; hand-tighten gland, then quarter-turn with strap wrench \\
\rowcolor{fogbg}
Erratic humidity readings & Condensation inside sensor bay & Inspect gasket seal; replace desiccant sachet (Part No.\ NX-DES-04) \\
Frequent low-power mode & Panel shaded or misaligned & Re-check southern exposure per Section~\ref{ch:install} \\
\bottomrule
\end{tabularx}
\end{center}

\begin{warnbox}
Never open the enclosure while the internal battery is under active solar charge in direct sunlight. Disconnect the panel lead at the MC4 connector first, or perform maintenance in shade / after dusk.
\end{warnbox}

\begin{tipbox}
Keep a spare desiccant sachet (Part No.\ NX-DES-04) in your field kit --- condensation is the single most common cause of erratic sensor readings reported to Nexa support.
\end{tipbox}

% ============================================================
\appendix
\chapter{Technical Specifications}
\label{app:specs}

\begin{center}
\renewcommand{\arraystretch}{1.25}
\begin{tabularx}{\textwidth}{@{}l X@{}}
\toprule
\rowcolor{steel}
\textcolor{white}{\bfseries Category} & \textcolor{white}{\bfseries Detail} \\
\midrule
Mechanical & Polycarbonate enclosure, IP66, UV-stabilized; 340 $\times$ 220 $\times$ 118\,mm; 2.4\,kg unpopulated \\
\rowcolor{fogbg}
Environmental & $-30\,^{\circ}\mathrm{C}$ to $+60\,^{\circ}\mathrm{C}$ operating; $-40\,^{\circ}\mathrm{C}$ to $+70\,^{\circ}\mathrm{C}$ storage; up to 100\% RH condensing \\
Power & 12\,V DC nominal input; 8\,W integrated monocrystalline solar panel; 18\,Ah LiFePO4 battery, field-replaceable \\
\rowcolor{fogbg}
Sensor bay & 6 $\times$ screw-terminal ports, hot-swappable, 0--5\,V and 4--20\,mA compatible \\
Communications & LTE-M/NB-IoT (3GPP Rel.\ 13); LoRaWAN 1.0.3 (EU868/US915); RS-485 optional add-on card \\
\rowcolor{fogbg}
Data retention & 30 days on-device buffer (flash), synced on reconnect after outage \\
Certifications & CE, FCC Part 15, IP66 (IEC 60529) \\
\rowcolor{fogbg}
Warranty & 3 years, hardware; battery consumable, 1 year \\
\bottomrule
\end{tabularx}
\end{center}

% ============================================================
\chapter{Quick Reference Card}
\label{app:quickref}

\noindent
\begin{minipage}[t]{0.48\textwidth}
\begin{tcolorbox}[colback=steel!6,colframe=steel,title=\textcolor{white}{\bfseries\faSignal\ \ Status LED},fonttitle=\bfseries,arc=1.5pt]
\small
\textbf{Green solid} --- normal, strong signal\\
\textbf{Green pulse} --- normal, weak signal\\
\textbf{Amber pulse} --- sensor fault\\
\textbf{Red solid} --- battery / radio failure
\end{tcolorbox}
\end{minipage}\hfill
\begin{minipage}[t]{0.48\textwidth}
\begin{tcolorbox}[colback=amberaccent!8,colframe=amberaccent,title=\bfseries\faBatteryFull\ \ Power Notes,fonttitle=\bfseries,arc=1.5pt]
\small
Low-power mode below 25\% charge\\
Full recharge: approx.\ 6 sunlight hours\\
Replace battery below 80\% rated capacity\\
Panel output check: $>$14\,V DC in daylight
\end{tcolorbox}
\end{minipage}

\par\vspace{0.5cm}

\noindent
\begin{minipage}[t]{0.48\textwidth}
\begin{tcolorbox}[colback=steel!6,colframe=steel,title=\textcolor{white}{\bfseries\faTools\ \ Torque Values},fonttitle=\bfseries,arc=1.5pt]
\small
Mounting bracket bolts: 4.5\,N\textperiodcentered m\\
Sensor compression glands: hand-tight $+$ 1/4 turn\\
Enclosure latch: hand-tight only, do not overtighten
\end{tcolorbox}
\end{minipage}\hfill
\begin{minipage}[t]{0.48\textwidth}
\begin{tcolorbox}[colback=amberaccent!8,colframe=amberaccent,title=\bfseries\faPhone\ \ Support,fonttitle=\bfseries,arc=1.5pt]
\small
Nexa Systems Field Support\\
Reference this unit's Device ID (enclosure label)\\
Have Nexa Cloud diagnostics export ready before calling
\end{tcolorbox}
\end{minipage}

\end{document}
